\documentclass{article}
\usepackage[a4paper, total={6in, 8in}]{geometry}

\usepackage[utf8]{inputenc}

\title{Data Analysis for Geoscience}
\author{Nathasya Christien}

\begin{document}

\maketitle

\section{Intro}
\subsection{Course syllabus}
\begin{itemize}
    \item Gridded dataset - ENSO index
    \item Trends - significange, temperature Trends
    \item Correlation and covariance - ENSO $\rightarrow$ rainfall, SST $\rightarrow$ rainfall
    \item Break-point analysis - define onset Indian monsoon
    \item FFT - temperature diurnal or annual [stationary analysis]
    \item Wavelets - African East Wind [non-stationary analysis]
    \item EOF - spatial, spatial rainfall, African monsoon onset - NAO
    \item Bonus: Linear stability analysis 
    \item Can be organized if there's interest: Climate feedback analysis
\end{itemize}

\subsection{Class remarks}
\begin{itemize}
    \item Command to print output and let us scroll: ncdump ... $\vert$ less 
    \item NC file in climate field follows climate format (CF) compliant
    \item $\vert$ is called pipe, tr is translate or replacing
    \item Always print (echo) our command first, especially when we use wild card (* or ?) or force deletation
    \item enlarge : choose one grid and project it
    \item cdo uses broadcasting method to match time axis. It's kind of like periodic extension for missing time index. Can be a problem when we have mismatch calender system.
    \item cdo mean will neglect missing values and only consider valid values as valid data; cdo avg will return missing value when it appears in our time series.
    \item If we do simple arithmetic average for global temperature, we will be biased towards the pole. The output would be too cold.
    \item Conditional commands with cdo to compare two files: ge, gt, eq, lt, le
    \item Conditional commands with cdo to compare a file with a constant: gec
    \item Writing file is more costly than some operators, so instead of writing the outputs, it's better to pipe the process (use - for cdo)
    \item There's an argument to tell cdo to instead parallelize process, but do he command sequentially
    \item panoply: software to plot well
    \item Always think about how to validate your processing or plotting result. It's either the output is wrong (there's something wrong with your methodology) or your current understanding of the expected value is wrong.
\end{itemize}

\subsection{General workflow example}
\begin{enumerate}
    \item Cutout region - sellonlatbox
    \item Merged in time - mergetime
    \item June to Sep - selmon
    \item Time average - timmean
    \item anomaly - yearmean sub
    \item Compute weighted area - fldmean
\end{enumerate}

\section{Trend Derivation}

\begin{itemize}
    \item Let us have sample of data $(x_i, y_i)$ for $i=1,...,n$. We can use these $n$ points to do linear regression. When we have a new point, the slope might change. In this class, we will deal with this certainty.
    \item Linear model: $y=mx+c$ where
    \begin{itemize}
        \item $y$ : dependant variable
        \item $x$ : independant variable
        \item $m,c$ : slope and offset
    \end{itemize}
    \item One metric is sum of squared errors (SSE): $\sum_i (y_i - (mx_i + c))^2$ where it has a global minimum (convex curve)
    \item By deriving SSE to $c$, we obtain $\hat{c} = \bar{y} - m\hat{x}$.
    \item Remember that $\sum_i(y_i - \bar{y}) = 0$ by identity.
    \item By deriving SSE to $m$ and substitute $\hat{c}$, we obtain
    \[ \hat{m} = \frac{\sum(x_i - \bar{x})(y_i - \bar{y})}{\sum(x_i - \bar{x})^2}. \]
    \item Note the numerator. If $x$ and $y$ are linearly dependant, we can see that the values end up to be something very positive or negative. If it's independant to each other (linearly), the values would approximate to zero.
    \item Uncertainty sources: Measurement error, nonlinearity (for certain range can be linear and non in the other)
    \item Read GLAM (general linear additive model) which can explain nonlinear relationship by extending predictor variables
    \item We derived $Var(\hat{m}) = \frac{\sigma^2}{\sum(x_i - \bar{x})^2}$ which explains why variability increases as we increase noise spread or decrease sample data.
    \item Standard error can be seen as the standard deviation of our slope estimator. It's influenced by residual variance, sample size, and spread of $x_i$. We care more about slope variance.
    \item We build conf. interval of the slope estimator using t-distribution for small samples and normal distribution for large samples. If the interval does not touch zero, we can say the slope is significance.
    \item In high latitude, the trend is stronger because of local feedback (melted ice, albedo changes) and heat transport change (in aquaplanet, it's significant towards the poles).
    \item The trend when we use 1901-2023 (stationary trend?) time series is less prominent than 1970-2023 (nonstationary trend?)
\end{itemize}

\section{Correlation and Covariance}
\begin{itemize}
    \item Larger noise gives larger covariance by trial.
    \item Covariance's unit is first variable's unit times second variable's unit. Correlation is dimensionless.
    \item Proportion Explained = $\frac{SS_{reg}}{SS_{tot}}$
    \item Notation: $s_x$ is sample standard deviation of $x$
    \item We proved that $r^2$ (squared of corellation) is the explained variability
    \item We proved $t$ statistics
    \[ t = r \sqrt{\frac{n-2}{1-r^2}} \]
    such that we have to be careful when we have large dataset because it will return a significant statistic
    \item We need to choose $n-2$ (degree of freedom) that is true. For example, by making autocorrelation of a variable, there's a time window to know its correlation. There's a limit for example after 5*12 months, the temperature time series is not corellated anymore. Thus, the number of independent sample points are actually lower.
\end{itemize}

\section{Breakpoint Analysis}
\begin{itemize}
    \item Offset change with high noise can be unnoticable. Instead, we can look by eye when we have slope change.
    \item For the Indian region $\partial w / \partial t$ is the local storage of water vapor change. Now, it's negilible. However, by future warming, it will actually be significant and explain the forecast of monsoon onset delay.
    \item Negative MFC = export of water vapor out the domain. We have descent (dry air dominates) - dry season, less thunderstorm
    \item When we have positive MFC = divergence and ascent. Because local storage is approximately zero, this translates to higher precipitation.
    \item Reanalysis data has problem because it's a push-pull interaction between model result being fit to observation. It can leave residuals for moisture budget, e.g model is wet but observation really dry
    \item By using segmentation analysis, we are proposing a robust definition for monsoon onset towards noise in our climate data or phenomena.
\end{itemize}

\section{Fourier Analysis}
\begin{itemize}
    \item By eye, we can easily identify the mode in a pattern, but we overlook the weak frequencies.
    \item FFT can be applied not only to a time series. We can do it in space or even multi-dimensional.
    \item For square wave, we can expect the first mode to be in phase with the square waves. The next mode will try to adjust the first mode to be flattened.
    \item Negative frequency can be noted as a clockwise rotation in the complex coordinate. If we're only analyzing a real time series (to find the mode, not the actual power or amplitude), we can neglect the negative frequencies. However, for bivariate analysis or a complex time series, we care for both positive and negative frequencies.
    \item For real data input, we can the symmetry between positive and negative frequencies so the imaginary part will be cancelled out.
    \item Convolution: See it as covariance between signal anomalies and a cosine function. If it's weakly correlated, the weight is close to zero which means that particular harmonic wave has a weak weight to the dataset.
    \item FFT treats our input by extending it to correlate it to a continuous and infinite time series. Thus, input with trend messes with the output. For example, for the first mode, it tries to capture the trend.
    \item Critical pre-process: 
    \begin{itemize}
        \item De-meaning is not really necessary.
        \item De-trend to ensure starting and ending point start at the same vertical point.
        \item Tapering to smooth both ends so we can deal with sharp gradients.
    \end{itemize}
    \item Hann window is great to deal with leakage (infinitely differential and smooths out edges well), but it's biased towards the middle of the time series. Thus, in climate research, we often use Tukey window that has wider window with 100\%.
    \item If we feed white noise (e.g numbers from Gaussian distribution) to FFT, we can expect equal power along the spectrum. Whereas for our example with temperature, we have red noise as it is stronger for low periods and and dissipates for higher periods.
    \item Exercises:
    \begin{itemize}
        \item Use other windows for the temperature example
        \item Put trend in the temperature example
        \item Apply FFT to ENSO (can we see 7 year cycle?)
        \item Remember that FFT is based on evenly spaced data, so we have to reinterpolate our data
    \end{itemize}
\end{itemize}

\section{Wavelets}
\begin{itemize}
    \item Fourier can't capture intermittency where signal can occur on certain period intervals.
    \item To solve this, we create a window that's focused on certain period - snippets of window time. We perform the same analysis where correlate our data with certain harmonic wave.
    \item It's like doing FFT but adding time dimension. If we take time average of our wavelet map, it's like FFT result.
    \item $\Omega_0$ spectrum: Using high frequency, makes it closer to the standard FFT because on each time index, they can have strong similarity. If it's small, we're just distorting the amplitude. It work well seperating analysis on each time index, but the frequency becomes wrong.
    \item The time scale of the wavelet wave length is slightly decreased from the corresponding harmonic wave. Think about the Gaussian envelope enabling a slope that lets peaks happened quicker than what we expect. For high frequency, the shift from the envelope affects less, while it's noticable for small $\omega_0$.
    \item We chose $\omega_0=6$ if we want to compare with literatures as it's the convention. However, as proved, $2\pi$ is better as we have better value for time scale near to 1, also the Fourier time close to 1.
    \item Usually our spacing for period bins, we use log intervals, such that we can give more bins to shorter waves.
    \item Again: Larger omega gives accurate picking frequency, but the time jump is worse. Small omega captures time jump better, but we have frequency leakage.
    \item Preprocessing data before wavelet: Detrend, remove mean, and take field mean.
    \item COI basically warns us that we're not performing full analysis by the boundaries of our time series, so we can be more cautious.
\end{itemize}

\section{Wavenumber Analysis}
\begin{itemize}
    \item We use \textbf{spectral leakage} to refer to space ringing, while \textbf{spectral ringing} is for time analysis.
    \item Leakage or ringing happens because again, FFT relied on periodicity, such that if our data is not periodic in time or space, we will have sudden jump in our idealized wave. New frequencies will be introduced by FFT result.
    \item Thus, we need to introduce tapering data too for our data in both time and space.
    \item Usually, we smooth out our 2D FFT result to remove white noise in the background and highlight the peaks. However, there has not been one convential method for this.
\end{itemize}

\end{document}